% ============================================================================
%  Annexe B — Données de mesure
%  Tableaux bruts ayant servi au chapitre 6. Permet la vérifiabilité.
% ============================================================================
\chapter{Données de mesure}
\label{ann:donnees}

\section{Relevé des exécutions}
\label{ann:releve-executions}

\aecrire{Tableau des exécutions retenues : identifiant, date, chaîne, issue,
durée. Produit par \texttt{scripts/collecte-mesures.sh}. Si le volume est
important, n'en présenter qu'un échantillon représentatif et indiquer
l'effectif total.}

\begin{table}[H]
\centering
\caption{Extrait du relevé des exécutions}
\label{tab:ann-executions}
\begin{tabular}{@{}llllr@{}}
\toprule
\textbf{Id.} & \textbf{Date} & \textbf{Chaîne} & \textbf{Issue} & \textbf{Durée (min)} \\
\midrule
 & & & & \\
\bottomrule
\end{tabular}
\end{table}

\section{Mesures de couverture}
\label{ann:releve-couverture}

\aecrire{Relevé des taux de couverture lignes, fonctions et branches sur la
période d'observation.}

\section{Relevé des défauts d'analyse statique}
\label{ann:releve-defauts}

\aecrire{Nombre de défauts par priorité aux dates de relevé.}

\section{Relevé des vulnérabilités de composants}
\label{ann:releve-vulns}

\aecrire{Nombre de composants et répartition des vulnérabilités par sévérité.}

\section{Protocole de mesure de reproductibilité}
\label{ann:protocole-repro}

\aecrire{Décrire pas à pas la procédure appliquée en
section~\ref{sec:eval-reproductibilite}, de manière qu'un tiers puisse la
rejouer.}
